\documentclass{article}
\usepackage[utf8]{inputenc}

\usepackage[backend=biber]{biblatex}
\addbibresource{bibliography.bib}

\usepackage
{amsmath,amssymb,hyperref,mathpartir,minted,simplebnf,
stmaryrd,tikz-cd}

\usepackage{amsthm}
\newtheorem{theorem}{Theorem}
\newtheorem*{theorem*}{Theorem}
\newtheorem{lemma}{Lemma}
\newtheorem{proposition}[lemma]{Proposition}
\newtheorem{corollary}{Corollary}
\theoremstyle{definition}
\newtheorem{definition}{Definition}
\theoremstyle{remark}
\newtheorem*{remark}{Remark}

\usepackage{newunicodechar}
\newunicodechar{Γ}{$\Gamma$}
\newunicodechar{Π}{$\Pi$}
\newunicodechar{λ}{$\lambda$}
\newunicodechar{ᵀ}{\textsuperscript{T}}
\newunicodechar{ᵁ}{\textsuperscript{U}}
\newunicodechar{ᵗ}{\textsuperscript{t}}
\newunicodechar{ᵘ}{\textsuperscript{u}}
\newunicodechar{ᵣ}{\textsubscript{r}}
\newunicodechar{⊔}{$\sqcup$}

\newcommand{\rk}[1]{\mathsf{rk}_{#1}}
\newcommand{\PTS}{\mathsf{PTS}_{rk}}
\newcommand{\Rt}{\widetilde{\mathsf R}}
\newcommand{\R}{\mathsf R}
\newcommand{\Good}{\mathsf{Good}}
\newcommand{\Case}{\mathsf{Case}}
\newcommand{\Inv}{\mathsf{Inv}}
\newcommand{\ev}{\operatorname{ev}}
\newcommand{\Add}{\mathsf{Add}}
\newcommand{\Mul}{\mathsf{Mul}}
\newcommand{\z}{\mathsf z}
\newcommand{\suc}{\mathsf s}
\newcommand{\nc}{\mathop{\neg}\nolimits_c}
\newcommand{\ac}{\mathbin{\wedge_c}}
\newcommand{\oc}{\mathbin{\vee_c}}
\newcommand{\iffc}{\mathbin{\leftrightarrow_c}}
\newcommand{\Up}{\mathsf{Up}}
\newcommand{\Da}{D_{\!a}}
\newcommand{\za}{\mathsf z_a}
\newcommand{\sa}{\mathsf s_a}
\newcommand{\adda}{\mathsf a_a}
\newcommand{\ma}{\mathsf m_a}
\newcommand{\Ga}{G_a}
\newcommand{\Ea}{E_a}
\newcommand{\Oa}{O_a}
\newcommand{\Gooda}{\mathsf{Good}_a}
\newcommand{\Casea}{\mathsf{Case}^{a}}
\newcommand{\Inva}{\mathsf{Inv}^{a}}
\newcommand{\nua}[1]{\nu^{a}_{#1}}

\title{A sound pure type system with ranks as sorts and
       Type-in-Type}
\author{
    Pavel Sokolov \\ HSE University, Laboratory of
    Theoretical Computer Science
    \and
    Alexander Kuklev \\ Radboud Univ. Nijmegen,
    Department of Software Science
}
\date{\today}

\begin{document}

\maketitle

\begin{abstract}
We present a dependently typed pure type system which
contains a circular axiom $s:s$ for a certain sort $s$
and is nonetheless strongly normalising. Girard's
paradox is avoided because our sorts are not closed
under a $\Pi$-former: function types are placed in
sorts according to their \textit{rank}. Turns out that
this system has a unique set of characteristics: it has
similar normalisation behaviour as STLC (so complexity
of normalisation is in \textsf{TOWER}) while at the
same time allowing non-uniform representation of
arbitrary elementary recursive functions; furthermore,
it has a nontrivial set-theoretical model in ZF+GCH.
Base normalisation proof is surprisingly simple (modulo
general metatheoretical properties of pure type
systems) and has been fully formalised in vanilla Rocq
(formerly Coq) without any additional axioms.
\end{abstract}

\section{Introduction}
Pure type systems (PTSs, introduced independently by
Berardi [CITE] and Terlouw [CITE]) are a family of
minimalistic typed $\lambda$-calculi including grammar
formers for variables $x\in\mathsf{Var}$, abstraction
$\lambda(x:T).u$, application $u\,t$, types of
dependent functions $\Pi(x:T).U$ (where $U$ might
depend on $x$) and \textit{sorts} $s\in\mathcal{S}$
which act like ``types of types''.

\begin{figure}[h]
\begin{bnf}
$x$ : \textsf{Variable} :in: $X$;;
$s$ : \textsf{Sort} :in: $S$;;
$T$, $U$, $t$, $u$ : \textsf{Term} ::=
| $x$ : variable
| $s$ : sort as type
| $\Pi(x\colon T).\,U$ : type of dependent functions
| $\lambda(x\colon T).\,u$ : anonymous function
| $u\,t$ : function application
;;
\end{bnf}
\end{figure}

Common syntax conventions to make it more lightweight
are the following: application associates to the left,
so $a\,b\,c$ is understood as $(a\,b)\,c$; in a
sequence of abstractions, all $\lambda$ symbols except
the first can be omitted; $T\to U$ is an alias for
$\Pi(x:T).U$ (where $U$ does not depend on $x$);
$(\to)$ is right-associative, so $T\to U\to V$ is
understood as $T\to(U\to V)$. In addition, even though
abstractions are always annotated with a type of an
argument, we are going to omit it on paper whenever
they are clear from the context. Also note that there
is no syntactical distinction between terms and types
--- they are defined by a single grammar; an expression
can act both as a term and as a type depending on a
side where it stands next to the colon in the typing
judgement $\Gamma\vdash\_:\_$.

A single PTS is determined by two relations:
\begin{itemize}
\item \textit{axiom} relation $\mathsf{Ax}\subseteq
    \mathcal{S}\times\mathcal{S}$ determines ``to which
    sort(s) a given sort belongs as an element'';
    instead of $(s,s')\in\mathsf{Ax}$ it is customary
    to write simply $s:s'$.
\item \textit{rule} relation $\mathsf{Rule}\subseteq
    \mathcal{S}^2\times\mathcal{S}$ places types of
    functions $\Pi(x:T).U$ into sorts accounting for
    the sorts where $T$ and $U$ are situated.
\end{itemize}

\begin{figure}[h]
\begin{mathpar}
\inferrule[Ax]{ }{\Gamma\vdash s:s'}
    (s, s')\in\mathsf{Ax}
\and
\inferrule[Rule]{
    \Gamma\vdash T:s_T \\
    \Gamma,x:T\vdash U:s_U
}{\Gamma\vdash\Pi(x:T).\,U:s}
(s_T, s_U, s)\in\mathsf{Rule}
\and
\inferrule[Var]{ }{\Gamma\vdash x:T}
    x:T\in\Gamma
\and
\inferrule[Abs]{
    \Gamma\vdash T:s_T \\
    \Gamma,x:T\vdash u:U
}{\Gamma\vdash\lambda(x:T).u:\Pi(x:T).U}
\and
\inferrule[App]{
    \Gamma\vdash u:\Pi(x:T).U \\
    \Gamma\vdash t:T
}{\Gamma\vdash u\,t:U[t/x]}
\and
\inferrule[Conv]{
    \Gamma\vdash t:T \\
    \Gamma\vdash U:s \\
    T =_\beta U
}{\Gamma\vdash t:U}
\end{mathpar}
\caption{Typing rules of pure type systems}
\label{typing}
\end{figure}

Systems from Barendregt's lambda cube [CITE] fit
naturally into the framework: for example, simply typed
lambda calculus (STLC) [CITE] can be viewed as a pure
type system with a single sort $*$, no axioms and a
single rule $(*,*,*)$. In general, however, pure type
systems are dependently typed since $\Pi$ allows the
codomain type to depend on the value of an argument ---
Coquand's Calculus of Constructions (CoC) [CITE] sits
comfortably at the top of lambda cube and is modeled as
a pure type system with two sorts $\{*,\square\}$, a
single axiom $*:\square$ and four rules corresponding
to simple functions $(*,*,*)$, polymorphic functions
$(\square,*,*)$, type operators
$(\square,\square,\square)$ and dependent types
$(*,\square,\square)$.

As is known [CITE] [CITE], both STLC and CoC are
strongly normalising and sound with respect to their
set-theoretical semantics. A classical result shows
that it is not true for arbitrary PTS.

\begin{theorem*}[Girard'72 [CITE]]
    The following PTS called System U is not
    normalising:
    \[\begin{array}{c}
        \mathcal{S}=\{*,\square,\triangle\},
        \quad\mathsf{Ax}=\{
            (*,\square),(\square,\triangle)
        \},
        \\ \mathsf{Rule}=\{
            (*,*,*),(\square,*,*),(\triangle,*,*),
            (\square,\square,\square),
            (\triangle,\square,\square)
        \}
    \end{array}\]
    Futhermore, the offending term has type
    $\Pi(P:*).P$.
\end{theorem*}

This result is dubbed ``Girard's paradox'' since it
renders System U inconsistent with respect to its
intended logical semantics: inhabitedness of
$\Pi(P:*).P$ would mean that every proposition is
provable. Thus any system where System U embeds is
inconsistent, too; in particular, an original version
[CITE] of Martin-L\"{o}f's Intuitionistic Type Theory
[CITE] embeds the following PTS known as
``Type-in-Type'' which is inconsistent due to Girard's
paradox:
\[\mathcal{S}=\{*\},\quad\mathsf{Ax}=\{*:*\},\quad
\mathsf{Rule}=\{(*,*,*)\}.\]
For a pure type system with usual set-theoretical
semantics, it is natural to expect that a sort of
\textit{small types} $*$ is closed under $\Pi$, thus we
cannot forbid the rule $(*,*,*)$; then, a circular
axiom $*:*$ has to go. Long story short, it motivates
an introduction of a \textit{universe hierarchy} [CITE]
$\mathcal{S}=\{\mathcal{U}_k\}_{k\in\mathbb{N}}$ with
axioms $\mathcal{U}_k:\mathcal{U}_{k+1}$ and rules
$(\mathcal{U}_k,\mathcal{U}_l,\mathcal{U}_{k\vee l})$
for each $k,l$; when interpreted in sets,
$\llbracket \mathcal{U}_k \rrbracket$ is taken to be
the $k$-th inaccessible cardinal [CITE].

Observing this, we took interest in the following two
questions:

\begin{itemize}
\item If we were to drop the rule $(*,*,*)$ and keep
    the axiom $*:*$ instead, can we find a nontrivial
    consistent pure type system?
\item Is there a consistent pure type system with total
    relations $\mathsf{Ax}$ and $\mathsf{Rule}$ whose
    natural set-theoretical semantics do not involve
    assumption of inaccessible cardinals?
\end{itemize}

Turns out that, for both questions, answer is ``yes''.
Furthermore, both yes-s are witnessed by a single PTS,
which we dub PTS-$rk$:
\[
\mathcal{S}=\{\mathsf{rk}_r\}_{r\in\mathbb{N}}
\quad
\mathsf{Ax}=
    \{(\mathsf{rk}_r,\mathsf{rk}_0)\}_{r\in\mathbb{N}}
\quad
\mathsf{Rule}=
    \{(\mathsf{rk}_r,\mathsf{rk}_s,
        \mathsf{rk}_{(r+1)\vee s})\}_{r,s\in\mathbb{N}}
\]
(Naming ``$\mathsf{rk}$'' comes from \textit{rank},
since rules compute precisely the rank of functions,
which is the maximal nesting of $\Pi$ in the domain
type.)

Indeed, there is a circular axiom
$\mathsf{rk}_0:\mathsf{rk}_0$, and System U is not
embeddable since $\mathsf{rk}_r$ is not closed
under $\Pi$ for any $r$. Also, as we will see in
Section~\ref{rk-sn}, our PTS-$rk$ is strongly
normalising despite an intuition that
$\mathsf{rk}_r:\mathsf{rk}_0$ might ``feel wrong''.
Modulo general metatheory of pure type systems laid out
in Section~\ref{pts-metatheory}, a normalisation proof
turned out to be suspiciously simple, even more so in
comparison with normalisation proofs of other
dependently typed systems [CITE]. To dispel our own
doubts, we formalised a proof of normalisation in Rocq
(formerly Coq) without using any additional axioms,
yielding a fully constructive proof; our journey in
formalisation is summarised in
Section~\ref{formalisation}.

Having completed the proof, we noticed that it exposes
a primitive recursive procedure to normalise terms of
PTS-$rk$; we dug further and discovered that, actually,
there is an intimate connection between a class of
elementary recursive functions [CITE] and our system,
as seen in Section~\ref{complexity}. As a byproduct of
our study on complexity, we built a sound set-theoretic
model for PTS-$rk$ (see Section~\ref{rk-model}) which
assumes \emph{Generalised Continuum Hypothesis} [CITE]
which is an unusual pick for dependently typed systems,
where \emph{inaccessible cardinals} are a more common
assumption [CITE]. We provide a closing summary and
some pointers on future directions of our research and
possible extensions of the system in
Section~\ref{conclusion}.

\subsection{Use of artificial intelligence}

At various points in our research, we interacted with
GPT-6 Astra [CITE] to prove something for us. Sometimes
it was a rehash of previously established results;
sometimes it presented something new. In case we used
the ideas hinted by LLM, we explicitly state so right
next to the proof; even so, the text is completely
human written since we wanted to make sure we
completely understood the results (in addition,
generated proofs always required additional
explanations or a complete rewrite to make them
understandable and simpler to read).

Also, we chatted with Claude Fable 5.1 to find better
points of reference.

\section{Metatheory of pure type systems}
\label{pts-metatheory}

For ease of presentation of our results, we assume a
Barendregt convention [CITE] that all bound variables
are distinct from free variables and from each other;
furthermore, $\alpha$-equivalence (renaming of bound
variables) is assumed as a trivial equality on terms.
Nevertheless, our formalisation takes variable
treatment seriously; see Section~\ref{formalisation}
for more details.

\subsection{PTS as a rewriting system}

In Fig.~\ref{typing}, a rule of \textit{conversion}
(\textsc{Conv}) mentions \textit{definitional equality}
of types $T=_\beta U$. It is defined as a reflexive,
symmetric, transitive closure of a one-step reduction.

\begin{definition}[One-step reduction]
Relation of a one-step reduction $t\to_\beta t'$ is the
smallest relation satisfying the following rules:
\begin{mathpar}
\inferrule{T \to_\beta T'}
    {\Pi(x:T).U \to_\beta \Pi(x:T').U}
\and
\inferrule{U \to_\beta U'}
    {\Pi(x:T).U \to_\beta \Pi(x:T).U'}
\and
\inferrule{T \to_\beta T'}
    {\lambda(x:T).t \to_\beta \lambda(x:T').t}
\and
\inferrule{t \to_\beta t'}
    {\lambda(x:T).t \to_\beta \lambda(x:T).t'}
\and
\inferrule{t \to_\beta t'}{t\,u \to_\beta t'\,u}
\and
\inferrule{u \to_\beta u'}{t\,u \to_\beta t\,u'}
\and
\inferrule[Beta]{}
    {(\lambda(x:T).t)\,u \to_\beta t[u/x]}
\end{mathpar}
\end{definition}

In essence, actual reduction happens only in rule
\textsc{Beta}; other rules simply state the we can
freely reduce any subterm.

The following result will be greatly useful in our
development. It was proven many times over for
different similar systems [CITE] [CITE]; we briefly
revisit the proof since previous presentations usually
did not allow reductions in the type of argument $T$ in
$\lambda(x:T).t$.

\begin{theorem}[Characterisation of $(=_\beta)$]
    In arbitrary PTS,
    \[
        t =_\beta u \quad\Leftrightarrow\quad
        \exists v.\;t\to_\beta^* v
        \land u\to_\beta^* v,
    \]
    where $(\to_\beta^*)$ is reflexive transitive
    closure of $(\to_\beta)$.
\end{theorem}
\begin{proof}
    Note that it is enough to prove that
    $(\to_\beta^*)$ has a diamond property --- for each
    $t_0,t_1,t_2$ such that $t_0\to_\beta^* t_1$ and
    $t_0\to_\beta^* t_2$ there exists a $t'$ such that
    $t_1\to_\beta^* t'$ and $t_2\to_\beta^* t'$:
    \[\begin{tikzcd}
        t_0 \arrow[r] \arrow[d]
        & t_1 \arrow[d, dashed] \\
        t_2 \arrow[r, dashed] & t'
    \end{tikzcd}\]
    Then, any ``mountains'' in the path $t=_\beta u$
    can be switched into ``valleys''.

    Now, diamond property is proven by standard
    technique of parallel reduction [CITE] which is
    defined as the smallest relation satisfying:
    \begin{mathpar}
    \inferrule{}{t\rightrightarrows t}
    \and
    \inferrule{
        T\rightrightarrows T' \\
        U\rightrightarrows U'
    }{\Pi(x:T).U\rightrightarrows\Pi(x:T').U'}
    \and
    \inferrule{
        T\rightrightarrows T' \\
        t\rightrightarrows t'
    }{\lambda(x:T).t\rightrightarrows\lambda(x:T').t'}
    \and
    \inferrule
        {t\rightrightarrows t'\\ u\rightrightarrows u'}
        {t\,u\rightrightarrows t'\,u'}
    \and
    \inferrule
        {t\rightrightarrows t'\\ u\rightrightarrows u'}
        {(\lambda(x:T).t)\,u\rightrightarrows t'[u'/x]}
    \end{mathpar}
    Intuitively, while $(\to_\beta)$ is a single
    $\beta$-reduction somewhere in the term and
    $(\to_\beta)^*$ is a sequence of reductions,
    $(\rightrightarrows)$ is several simultaneous
    reductions immediately available in the term.
    Notice, then, that $(\to_\beta)\subseteq
    (\rightrightarrows)\subseteq(\to_\beta^*)$; by
    monotonicity of reflexive transitive closure,
    $(\to_\beta^*)\subseteq
    (\rightrightarrows^*)\subseteq(\to_\beta^*)^*
    =(\to_\beta^*)$, so
    $(\rightrightarrows^*)=(\to_\beta^*)$. Thus it is
    enough to prove diamond property of
    $(\rightrightarrows^*)$. Diamond property is
    preserved under reflexive transitive closure, so it
    is enough to prove diamond propery of
    $(\rightrightarrows)$, which holds by tedious case
    analysis.
\end{proof}

Setting aside practical applications of this theorem
(to algorithmically check equivalence of terms, we can
now search for common reduction result), it will prove
useful already in our proof development. For example,
the following ``equality inversion'' lemmas are direct
corollaries:

\begin{corollary}[Equality inversion]
    If $s=_\beta s'$, then $s=s'\in\mathcal{S}$; if
    $\Pi(x:T).U=_\beta\Pi(x:T').U'$, then $T=_\beta T'$
    and $U=_\beta U'$.
\end{corollary}

They will become necessary in proving subject
reduction, or type preservation. But first, we will
show some structural properties of a typing judgement
given in Fig.~\ref{typing}.

\subsection{Properties of a typing judgement}

Again, they surely have been proven before for PTS
[CITE], but we slightly differ in the formulation of
typing rules in important places so it is better to
recheck that these hold.

\begin{lemma}[Weakening]
    The following rule is admissible:
    \[\dfrac{\Gamma,\Delta\vdash t:T}
            {\Gamma,x:U,\Delta\vdash t:T}.\]
\end{lemma}
\begin{proof}
    By structural induction on the typing derivation.
\end{proof}

\begin{lemma}[Substitution]
    The following rule is admissible:
    \[\dfrac{
        \Gamma,x:U,\Delta\vdash t:T\quad
        \Gamma\vdash u:U
    }{\Gamma,\Delta[u/x]\vdash t[u/x]:T[u/x]}.\]
\end{lemma}
\begin{proof}
    By structural induction on the typing derivation
    for $t$. For variables in $\Delta$, it is necessary
    that both $\Delta$ and $T$ get substitutions of $u$
    for $x$ so \textsc{Var} is still applicable; for
    places in $t$ where $x$ used to be we need to apply
    weakening lemma to the typing derivation for $u$
    exactly $|\Delta|$ times. In the case of
    \textsc{Conv} it is important that definitional
    equality of types is preserved by substitution
    (thanks to characterisation of $(=_\beta)$).
\end{proof}

In addition, a family of so-called \textit{inversion
lemmas} [CITE] stating that type of well-typed term is
largely determined by shape of a term are also proven
by straightforward induction on typing derivation (each
one is proven separately, no mutual induction
involved):

\begin{lemma}[Inversion lemmas]
\begin{align}
    \Gamma\vdash s:T &\Rightarrow
    \exists s'. (s,s')\in\mathsf{Ax} \land T =_\beta s'
    \\
    \Gamma\vdash\Pi(x:T).U:V &\Rightarrow
    \exists s_T,s_U,s.\left\{\begin{array}{l}
        (s_T,s_U,s)\in\mathsf{Rule} \\
        \Gamma\vdash T:s_T \\
        \Gamma,x:T\vdash U:s_U \\
        V =_\beta s
    \end{array}\right. \\
    \Gamma,x:T,\Delta\vdash x:U &\Rightarrow T=_\beta U
    \\
    \Gamma\vdash\lambda(x:T).t:V &\Rightarrow
    \exists s,U.\left\{\begin{array}{l}
        \Gamma\vdash T:s \\
        \Gamma,x:T\vdash t:U \\
        V =_\beta \Pi(x:T).U
    \end{array}\right. \\
    \Gamma\vdash t\,u:V &\Rightarrow
    \exists U,T.\left\{\begin{array}{l}
        \Gamma\vdash t:\Pi(x:U).T \\
        \Gamma\vdash u:U \\
        V =_\beta T[u/x]
    \end{array}\right.
\end{align}
\end{lemma}

To progress further, we need to make certain
assumptions on the PTS we have in question.

\begin{definition}[Functional PTS]
    A pure type system is called \emph{functional} if
    $\mathsf{Ax}\subseteq\mathcal{S}\times\mathcal{S}$
    and $\mathsf{Rule}\subseteq\mathcal{S}^2\times
    \mathcal{S}$ are functional relations, that is,
    \[
        (s:s')\land(s:s'')\Rightarrow s'=s''
        \quad\text{and}\quad
        (s_1,s_2,s)\in\mathsf{Rule}\land(s_1,s_2,s')\in
        \mathsf{Rule}\Rightarrow s=s'
    \]
\end{definition}

\begin{definition}[Total PTS]
    A pure type system is called \emph{total} if
    $\mathsf{Ax}\subseteq\mathcal{S}\times\mathcal{S}$
    and $\mathsf{Rule}\subseteq\mathcal{S}^2\times
    \mathcal{S}$ are total relations, that is,
    \[
        \forall s.\exists s'.(s,s')\in\mathsf{Ax}
        \quad\text{and}\quad
        \forall s_1,s_2.\exists s.(s_1,s_2,s)\in
        \mathsf{Rule}
    \]
\end{definition}

\begin{theorem}[Uniqueness of typing]
    In a functional PTS, the following rule is
    admissible:
    \[\dfrac{\Gamma\vdash t:T \quad \Gamma\vdash t:T'}
        {T =_\beta T'}.\]
\end{theorem}
\begin{proof}
    By structural induction on $t$ and
    inversion of given typing derivations.
    \begin{itemize}
    \item $t=s\in\mathcal{S}$: by functionality of
        $\mathsf{Ax}$.
    \item $t=\Pi(x:T).U$: inversion of typing gives a
        pair of (possibly different) sorts for $T$ and
        another pair for $U$. By induction we get that
        ranks inside pairs are definitionally equal,
        which by equality inversion gives strictly
        equal sorts. Then what is needed holds by
        functionality of $\mathsf{Rule}$.
    \item Case for $t=x\in\mathsf{Var}$ holds
        trivially.
    \item $t=(\lambda(x:T).u)$: by induction hypotheses
        and by congruence of $(=_\beta)$.
    \item $t=u\,v$: by typing inversion and by
        induction hypothesis for $u$ we get $u:
        \Pi(x:T_1).U_1$ and $\Pi(x:T_1).U_1=_\beta
        \Pi(x:T_2).U_2$. We need $U_1[v/x]=_\beta
        U_2[v/x]$ which we get from equality inversion
        and preservation of def.eq. under substitution.
    \end{itemize}
\end{proof}

In our proof of normalisation, we are going to rely on
yet another property: if there is a derivation for
$t:T$, then $T$ is itself typeable with some sort.
However, it does not hold in our current formulation of
typing rules: in $\Gamma\vdash x:T$, nothing prevents a
context $\Gamma$ to contain an entry $(x:T')$ where
$T'$ is not typeable. Thus we need an additional
assumption.

\begin{definition}[Context well-formedness]
    Context $\Gamma$ is \emph{well-formed} (written
    $\vdash\Gamma$) if every type in it is typed by
    some sort. Formally, set of well-formed contexts is
    the least set closed by the following rules:
    \[
        \dfrac{ }{\vdash\cdot} \quad\text{and}\quad
        \dfrac{\vdash\Gamma\quad\Gamma\vdash T:s}
              {\vdash\Gamma,x:T}
    \]
    (where $\cdot$ denotes an empty context.)
\end{definition}

Usually, when presenting a dependently typed system,
context well-formedness is added as a premise to typing
rules which would otherwise be premise-free (such as
\textsc{Ax} and \textsc{Var} in a PTS). Here we instead
provide it as an additional assumption only when
necessary; in addition to leading to simpler metatheory
(we avoid an otherwise mutually inductive definition of
$(\vdash\_)$ and $(\_\vdash\_:\_)$), this is closer to
how decision procedures for typing judgements are
implemented.

\begin{theorem}[Typing yields correct type]
\label{type-is-correct}
    In a total PTS, if $\Gamma\vdash t:T$ and
    $\vdash\Gamma$, then there exists $s\in\mathcal{S}$
    such that $\Gamma\vdash T:s$.
\end{theorem}
\begin{proof}
    By structural induction on typing derivation.
    \begin{itemize}
    \item \textsc{Ax}, \textsc{Rule}: by totality of
        $\mathsf{Ax}$.
    \item \textsc{Var}: follows from $\vdash\Gamma$.
    \item \textsc{Abs}: by totality of $\mathsf{Rule}$.
    \item \textsc{App}: $\Gamma\vdash t\,u:T[u/x]$.
        From induction hypothesis for $t$ we know that
        $\Gamma\vdash t:\Pi(x:U).T:s$ for some $s$,
        then by typing inversion for $\Pi(x:U).T$ we
        get that $T$ is typable with some sort. Then
        what is needed follows from substitution lemma.
    \item \textsc{Conv}: what is needed is present in
        the premise.
    \end{itemize}
\end{proof}

Totality is not a strictly necessary requirement ---
similar true propositions can be stated for non-total
PTS, for example in CoC typing yields type which is
either itself typable or equal to $\square$ [CITE]. But
PTS-$rk$ is indeed total so the theorem above is enough
for our purposes.

\begin{theorem}[Subject reduction]
    In a total PTS, the following rule is admissible:
    \[\dfrac{
        \Gamma\vdash t:T \quad \vdash\Gamma
        \quad t \to_\beta t'
    }{\Gamma\vdash t':T}\]
\end{theorem}
\begin{proof}
    By structural induction over derivation of
    $t\to_\beta t'$. Inversion of a typing judgement
    yields all necessary info, albeit because of it
    types are now given only up to definitional
    equality so we have to apply conversion rule in
    each case. Crucially, it requires the type we
    convert into to be typable by a sort, so that is
    where we use Theorem~\ref{type-is-correct}.

    The only tricky case are reductions of a domain
    type in $\Pi(x:T).U$ and $\lambda(x:T).t$. Let us
    consider the case of the following reduction:
    \[\lambda(x:T).t \to_\beta \lambda(x:T').t.\]
    In order to complete this case, we have to prove
    that $\Gamma,x:T'\vdash t:U$ given that
    $\Gamma,x:T\vdash t:U$. By appropriately
    strengthening the statement, we get that it is
    enough to prove admissibility of the following rule
    called \emph{context reduction}:
    \[\dfrac{
        \Gamma,x:U,\Delta\vdash t:T\quad
        \vdash\Gamma,x:U,\Delta\quad U\to_\beta U'
    }{\Gamma,x:U',\Delta\vdash t:T}\]
    The proof is, again, by induction on the typing
    derivation. The only non-trivial case is with
    \textsc{Var} rule used on $x$; it is taken care of
    by appropriate use of \textsc{Conv} rule to convert
    $x:U'$ to $x:U$. To use conversion, we have to
    prove that $U$ is well-typed; thankfully, it
    follows from $\vdash\Gamma,x:U,\Delta$.

    Note that context reduction is admissible in
    arbitrary PTS! We do not need it anywhere else,
    though, so we did not state it as a separate lemma.
\end{proof}

Characterisation of definitional equality, uniqueness
of typing and subject reduction together imply the
following corollary we are going to use a few times in
our normalisation proof:

\begin{corollary}[Type stability]
    In a total functional PTS, the following holds:
    \[\left\{\begin{array}{l}
        \Gamma\vdash t:T\\
        \Gamma\vdash t':T'\\
        \vdash\Gamma\\
        t =_\beta t'
    \end{array}\right.\quad\Rightarrow\quad
    T=_\beta T'.\]
\end{corollary}

\subsection{Properties of strongly normalising terms}

We are almost done with classical results! What is left
now is to state which terms are strongly normalising
and collect properties of s.n. terms which allow our
soundness proof to go through.

Informally, a term is strongly normalising iff there
are no infinite reduction chains starting from it.
In constructive setting and in formalisation, the
following definition is a bit better [CITE]:

\begin{definition}[Strongly normalising term]
\label{sn-definition}
    Set of s.n. terms is the smallest set $\mathsf{sn}$
    in which the following rule is valid:
    \[\dfrac{\forall t',(t\to_\beta t')\Rightarrow
        t'\in\mathsf{sn}}{t\in\mathsf{sn}}.\]
    In other words, a term is strongly normalising iff
    all possible reductions from it are strongly
    normalising.
\end{definition}

Thus, terms which cannot be reduced further (also known
as terms \emph{in normal form}) are strongly
normalising; then, terms which directly reduce into
normal forms are also strongly normalising; and so on.
Therefore, by this definition, strongly normalising
terms are those that every reduction chain starting
from them is finite.

The following properties are then provable by
definition of $\mathsf{sn}$ and by structural induction
on available derivations of $\mathsf{sn}$ for terms in
question:

\begin{lemma}[Structural properties of $\mathsf{sn}$]
\label{sn-props}
\begin{mathpar}
    \inferrule{}{x,s\in\mathsf{sn}}\and
    \inferrule{T\in\mathsf{sn}\\ U\in\mathsf{sn}}
        {\Pi(x:T).U\in\mathsf{sn}}\and
    \inferrule{T\in\mathsf{sn}\\ t\in\mathsf{sn}}
        {\lambda(x:T).t\in\mathsf{sn}}\and
    \inferrule{t\,u\in\mathsf{sn}}{t\in\mathsf{sn}}\and
    \inferrule{t[u/x]\in\mathsf{sn}}{t\in\mathsf{sn}}
\end{mathpar}
\end{lemma}

Application, of course, is the only nontrivial case. In
case of $\beta$-reduction, strong normalisation of both
$(\lambda(x:T).t)$ and $u$ does not imply anything
about normalisation behaviour of $t[u/x]$: consider a
classical example from untyped lambda calculus [CITE],
$\omega=(\lambda(x:T).x\,x)$ and
$\Omega=\omega\,\omega$; while obviously not typable,
$\omega$ is strongly normalising, while $\Omega$
decidedly is not.

Situation changes when a head of application is not an
abstraction at all. Thus we arrive at the definition of
an \textit{atomic term}.

\begin{definition}[Atomic term]
    Set of atomic terms is the smallest set satisfying
    the following conditions:
    \begin{enumerate}
        \item $s$, $\Pi(x:S).T$ and $y$ are atomic for
            all variables $x$,$y$, sort $s$ and types
            $S$,$T$;
        \item If $t$ is atomic, then $t\,u$ is also
            atomic.
    \end{enumerate}
\end{definition}

Atomic terms are useful precisely because their strong
normalisation survives application:

\begin{lemma}
\label{sn-atomic}
    If $t,u\in\mathsf{sn}$ and $t$ is atomic, then
    $t\,u\in\mathsf{sn}$.
\end{lemma}
\begin{proof}
    By induction on derivations of $\mathsf{sn}$ for
    $t$ and $u$. Consider possible redexes in $t\,u$:
    $t\,u$ is not itself a redex since $t$ is atomic
    and hence not a lambda; the only possible redexes
    are either inside $t$ or inside $u$. Applying
    corresponding induction hypotheses, we reach the
    desired result. The only caveat is that, for
    induction to go through, we need to check that
    atomicity of $t$ is preserved by reduction; this is
    indeed true by induction on derivation of atomicity
    for $t$.
\end{proof}

A dual view on strong normalisation of application is
the following: if we already know that
$t[u/x]\in\mathsf{sn}$, then the most tricky case is
already taken care of.

\begin{lemma}[Weak head expansion [CITE]]
\label{sn-whe}
    \[\left\{\begin{array}{l}
        t[u/x]u_1\ldots u_n\in\mathsf{sn} \\
        U,u\in\mathsf{sn}
    \end{array}\right.\Rightarrow
    (\lambda(x:U).t)u u_1\ldots u_n\in\mathsf{sn}.\]
\end{lemma}
\begin{proof}
    By induction on $n$.
    \begin{itemize}
        \item $n=0$. We need to prove that if
            $t[u/x],u,U\in\mathsf{sn}$, then
            $(\lambda(x:U).t)u\in\mathsf{sn}$.
            Note that $t[u/x]\in\mathsf{sn}$ implies
            $t\in\mathsf{sn}$; then, by structural
            induction on derivations of $\mathsf{sn}$
            for $t,u,U$ we get what we want by
            considering all possible reductions.
        \item $n\to(n+1)$. By induction hypothesis, it
            is enough to show the following:
            \[\left\{\begin{array}{l}
                (\lambda(x:U).t)u u_1\ldots u_n
                    \in\mathsf{sn} \\
                t[u/x]u_1\ldots u_{n+1}\in\mathsf{sn}
            \end{array}\right.\Rightarrow
                (\lambda(x:U).t)u u_1\ldots u_n u_{n+1}
                    \in\mathsf{sn}.
            \]
            This is true by structural induction on
            available evidence.
    \end{itemize}
\end{proof}

\section{Strong normalisation of PTS-$rk$}
\label{rk-sn}

\begin{theorem}[Strong normalisation of PTS-$rk$]
\label{normalisation}
    \[ \Gamma\vdash t:T\,\land\vdash\Gamma
    \Rightarrow t \in\mathsf{sn}. \]
\end{theorem}
\begin{proof}
    By induction on the typing derivation. Case
    \textsc{Conv} is trivial; due to
    Lemma~\ref{sn-props}, cases \textsc{Ax},
    \textsc{Rule}, \textsc{Var} and \textsc{Abs}
    present no problem as well. The only problematic
    case is application:
    \[
        \left\{\begin{array}{l}
            \Gamma\vdash u:\Pi(x:T).U \\
            \Gamma\vdash t:T \\
            \vdash\Gamma \\
            t,u\in\mathsf{sn}
        \end{array}\right.\Rightarrow
        u\,t\in\mathsf{sn}
    \]
    Now, here is the part where stratification
    according to ranks comes in useful. As typing
    yields correct type, there is some $r$ such that
    $\Gamma\vdash \Pi(x:T).U:\mathsf{rk}_r$. Let us
    explicitly add this to the list of premises; we
    will call the following proposition
    \emph{Application Lemma}:
    \[ \mathrm{AL}_r \triangleq \quad
        \text{``}\left\{\begin{array}{l}
            \Gamma\vdash \Pi(x:T).U:\mathsf{rk}_r\\
            \Gamma\vdash u:\Pi(x:T).U \\
            \Gamma\vdash t:T \\
            \vdash\Gamma \\
            t,u\in\mathsf{sn}
        \end{array}\right.\Rightarrow
        u\,t\in\mathsf{sn}\text{''}
    \]
    What indexing by $r$ bought us is that now we can
    try proving this lemma by induction on $r$. Again,
    as it always is with strong normalisation of
    application, the problem is in case where
    application head is a lambda and a
    $\beta$-reduction takes place. This calls for
    another lemma, \emph{Substitution Lemma}:
    \[ \mathrm{SL}_r \triangleq \quad
        \text{``}\left\{\begin{array}{l}
            \Gamma,x:T,\Delta\vdash u:U\\
            \Gamma\vdash T:\mathsf{rk}_r\\
            \Gamma\vdash t:T\\
            \vdash\Gamma\\
            t,u\in\mathsf{sn}
        \end{array}\right.\Rightarrow
        u[t/x]\in\mathsf{sn}\text{''}
    \]
    A proof of $\mathrm{AL}_r$ for all $r$ (which is
    what we are after) is done by complete mutual
    induction with $\mathrm{SL}_r$. Thus there are two
    main parts:
    \begin{itemize}
        \item $(\forall s<r.\;\mathrm{SL}_s)
            \Rightarrow \mathrm{AL}_r$: by induction on
            derivations of $u,t\in\mathsf{sn}$.
            Consider possible reductions of $u\,t$:
            most of them happen either in $u$ or in
            $t$, so that results of reductions from
            $u\,t$ are strongly normalising by
            inductive hypotheses. The only problematic
            case is when $u=\lambda(x:T').v$ and the
            reduction is $\beta$-reduction. By
            inversion of typing for $u$, it should be
            the case that $T'=_\beta T$; and by
            inversion of typing for $\Pi(x:T).U$, we
            get that $T:\mathsf{rk}_s$ and
            $U:\mathsf{rk}_t$ such that
            $r=\max(s+1, t)$. Then, by type stability,
            $T':\mathsf{rk}_s$, and furthermore, $s<r$,
            so we can apply substitution lemma.
        \item $(\forall s\le r.\;\mathrm{AL}_s)
            \Rightarrow \mathrm{SL}_r$. So we need to
            prove $u[t/x]\in\mathsf{sn}$. We do it by
            simultaneous induction on derivation of
            $u\in\mathsf{sn}$ and structural induction
            on $u$. As substitution propagates through
            $u$, most cases are satisfied by
            Lemma~\ref{sn-props}. A funny case is when
            $u=x$: then, a premise that
            $t\in\mathsf{sn}$ helps.

            The only problematic case is
            application $u=u_1\,u_2$; we consider three
            cases here:
            \begin{itemize}
                \item $u_1[t/x]$ is atomic --- handled
                    by Lemma~\ref{sn-props}.
                \item $u_1$ is not atomic --- by virtue
                    of it being well-typed, it means
                    that $u=(\lambda(x:V).u')v_1
                    \ldots v_k u_2$. All possible
                    reductions from $u[t/x]$ are then
                    handled by induction hypotheses.
                \item $u_1$ is atomic, but $u_1[t/x]$
                    is not --- by virtue of $u_1$ being
                    well-typed, the only case when
                    this is possible is when
                    $u_1 = x v_1\ldots v_k$, thus
                    $u[t/x]=t (v_1[t/x])\ldots
                    (v_k[t/x])(u_2[t/x])$. From
                    induction hypotheses we know that
                    all $v_i[t/x]$ and $u_2[t/x]$ are
                    s.n.; $t$ is s.n. by premise and
                    its rank is $r$. Then $t v_1[t/x]$
                    is s.n. by application lemma for
                    $r$; but note that rank of
                    $t v_1[t/x]$ is not greater than
                    $r$ so application lemma applies
                    again. Thus, by induction on $k$,
                    the whole chain of applications is
                    s.n.
            \end{itemize}
    \end{itemize}
\end{proof}

This is a retelling of a proof that our friend,
Aleksandr Levin, got after half-jokingly inputting the
following prompt into GPT-6 Pro:

\begin{quote}
\begin{verbatim}
Consider the following pure type system:

S = { rk_n | n natural }
Axiom = { (rk_n, rk_0) | n natural }
Rule = { (rk_m, rk_n, rk_max(m+1,n)) | m, n natural }

Is it strongly normalizing or not?
\end{verbatim}
\end{quote}

After approximately 20 minutes, LLM outputted the proof
which we retold above. Looking at it in retrospect, it
seems as if we stratified functions by ranks with
intent for such a direct proof to go through; honestly,
we did not think of that and were eagerly trying to
adapt well-known approaches for proving normalisation
of pure type systems employing either saturated sets or
candidates of reducibility [CITE] [CITE], which did not
work because our ``hierarchy'' of ranks is not
well-founded enough because of axioms
$\mathsf{rk}_n:\mathsf{rk}_0$. This proof, however,
does not even bother with it which means that we could
pick \emph{any} relation as $\mathsf{Ax}$.
Nevertheless, putting all sorts into $\mathsf{rk}_0$
seems the most sensible approach to us since,
informally, sort is kind of a constant which cannot be
converted into a function.

\subsection
{Formalising normalisation of PTS-$rk$ in Rocq}
\label{formalisation}

In view of high technicality of the proof for strong
normalisation which involves wrangling of unusual
definition for strong normalisation and complex
induction patterns (and in view of the fact that proof
was supplied by the probabilistic machine), we decided
to formalise it so that a sound simple algorithm would
check it for us. One of the natural choices for
formalisation of programming languages is Rocq
(formerly Coq) [CITE], an interactive theorem prover
based on Polymorphic, Cumulative Calculus of Inductive
Constructions [CITE]. As is the case with similar
typed calculi, it is both a programming language and a
higher-order logic, and correctness of a proof is
checked by the type checker; in addition to that, Rocq
features powerful metaprogramming facilities which
allow writing custom \emph{tactics} --- programs
written in a meta-language which generate terms of the
target language, greatly simplifying the proofs. While
set-theoretical semantics of Rocq assume an existence
of countable hierarchy of inaccessible cardinals
[CITE], our formalisation never delves into the
hierarchy and can be formalised in a much more
restricted theory.

In order to formalise our strong normalisation proof,
we first need to provide a general metatheory of PTS as
laid out in Section~\ref{pts-metatheory}. We could
reuse existing projects on PTS formalisations such as
[CITE] but we wanted to explore the topic ourselves
which led to a few minor differences discussed
previously, thus we started from scratch.

When one formalises syntax with binders in a proof
assistant, one of the most important decisions is how
to represent binders. Here we went with making our
terms indexed by the number of free variables it uses,
a standard technique originating from [CITE]:

\begin{minted}{coq}
Inductive Term (L: Type) (n: nat) : Type :=
| sort: L -> Term L n
| Π: Term L n -> Term L (S n) -> Term L n
| var: Fin n -> Term L n
| lam: Term L n -> Term L (S n) -> Term L n
| app: Term L n -> Term L n -> Term L n.
\end{minted}

Here, $L$ is a type of sorts ($L$ comes from
\emph{level} which is a name used for it in universe
hierarchies), and $n$ is a number of free variables.
Entering new scope (such as in $\Pi$ and in
$\mathtt{lam}$) adds one new free variable; variables,
then, have to be numbers in range $[0;n)$, which are
usually represented by a $\mathtt{Fin}$ type.

What is convenient in this representation is that it
makes it harder to make mistakes when wrangling with
context, defining substitution etc. --- all of the
techinical details which make working with syntax with
binders so hard (even de Bruijn made mistakes in
METAMATH when first implementing it! [CITE]). Also it
has a certain category-theoretical feel to it:
sufficiently generalised operations of renaming and
substitution give the typal representation of syntax a
structure of a \emph{relative monad} [CITE]. It is
really hard to go wrong when implementing these
operations (and the natural implementation turns
variables into de Bruijn indices [CITE]); on top of
them, usual substitution is easily definable.

Most of the proof development in Rocq mirrors what is
written in Section~\ref{pts-metatheory}; in fact, we
first wrote a proof in Rocq and then retold it in
written English here, omitting tedious technicalities.
The only considerable difference is in handling of the
set of strongly normalising terms since, obviously, the
proof of Substitution lemma given above does not
follow the structure of inductive definition
\ref{sn-definition}. Instead, we use a syntax-directed
characterisation of the set of strongly normalising
terms introduced in \cite{Joachimski2003Short}:

\begin{definition}[Syntax-directed characterisation
    of strong normalisation]
    Sets of terms $\mathsf{SN}$ (strongly normalising),
    $\mathsf{SNx}$ (strongly normalising neutral) and a
    relation $\to_\mathsf{WHR}$ (weak head reduction)
    are the least sets where the following rules are
    valid:
    \begin{mathpar}
        \inferrule{ }{\mathsf{SNx}\,s}(s\in\mathcal{S})
        \and
        \inferrule{\mathsf{SN}\,T\\\mathsf{SN}\,U}
                  {\mathsf{SNx}(\Pi(x:T).U)}
        \and
        \inferrule{ }{\mathsf{SNx}\,x}
        \and
        \inferrule{\mathsf{SNx}\,t\\\mathsf{SN}\,u}
                  {\mathsf{SNx}\,(t\,u)}
        \\
        \inferrule{\mathsf{SNx}\,t}{\mathsf{SN}\,t}
        \and
        \inferrule{\mathsf{SN}\,T\\\mathsf{SN}\,t}
                  {\mathsf{SN}\,(\lambda(x:T).t)}
        \and
        \inferrule{t\to_\mathsf{WHR}u\\\mathsf{SN}\,u}
                  {\mathsf{SN}\,t}
        \and
        \inferrule{\mathsf{SN}\,U\\\mathsf{SN}\,u}
                  {(\lambda(x:U).t)\,u\to_\mathsf{WHR}
                  t[u/x]}
        \and
        \inferrule{t\to_\mathsf{WHR}t'}
                  {t\,u\to_\mathsf{WHR}t'\,u}
    \end{mathpar}
\end{definition}

In a sense, $\mathsf{SNx}$ is a description of atomic
strongly normalising terms; $\mathsf{SN}$ adds a case
for lambdas and for weak head expansion, covering all
possible cases. Existence of such sets is by virtue of
them being subsets of the set of terms, so such mutual
inductive definition is not problematic.

One can prove that set $\mathsf{SN}$ is equal to
$\mathsf{sn}$; we, however, need only the following:

\begin{lemma}[Soundness of $\mathsf{SN}$]
    $\mathsf{SN}\subseteq\mathsf{sn}$.
\end{lemma}
\begin{proof}
    By mutual induction on the evidence
    $\mathsf{SN}\,t$ for arbitrary $t$. To do this, we
    also have to provide a term-indexed proposition we
    are proving for $\mathsf{SNx}$ and a relation on
    terms we are proving for $\mathsf{WHR}$:
    \begin{itemize}
        \item $\forall t,\mathsf{SNx}\,t\Rightarrow
            \mathsf{sn}\,t$.
        \item $\forall t\,t',\,(t\to_\mathsf{WHR}t')
            \Rightarrow(t\to_\mathsf{WHR-sn}t')$ (where
            the last relation is defined in the same
            way as WHR, but the head contains evidence
            of $\mathsf{sn}$ instead of $\mathsf{SN}$.)
    \end{itemize}
    Proof is extremely simple: by
    Lemmas~\ref{sn-props},\ref{sn-atomic},\ref{sn-whe}.
    We just have to notice that if term is
    $\mathsf{SNx}$, it is atomic, so lemma
    \ref{sn-atomic} applies in the relevant case.
\end{proof}

Then, to prove strong normalisation, it is enough to
prove that if term is typable, it is $\mathsf{SN}$; we
only have to state Application Lemma and Substitution
Lemma in terms of $\mathsf{SN}$, and the proof goes
through.

\begin{lemma}
    Application Lemma for all $s\le r$ implies
    Substitution Lemma for $r$.
\end{lemma}
\begin{proof}
    We have to prove that $\mathsf{SN}\,t$ and
    $\mathsf{SN}\,u$ imply $\mathsf{SN}\,(t[u/x])$.
    We do so by mutual induction on evidence for $t$
    with the following propositions proven for
    $\mathsf{SNx}$ and WHR:
    \begin{itemize}
        \item $\mathsf{SNx}\,t\land\Gamma\vdash t:T
            \Rightarrow\mathsf{SNx}\,(t[u/x])\lor
            (\mathsf{SN}\,(t[u/x])\land\exists s\le r.
            \;\Gamma\vdash T:\mathsf{rk}_s)$;
        \item $(t\to_\mathsf{WHR}t')\Rightarrow
            (t[u/x]\to_\mathsf{WHR}t'[u/x])\land
            (\mathsf{SN}\,(t'[u/x])\Rightarrow
            \mathsf{SN}\,(t[u/x]))$.
    \end{itemize}
    The proof then follows as written in a non-formal
    treatment above.
\end{proof}

\section{Computational complexity of PTS-$rk$}
\label{complexity}

(Proofs in the following sections are not formalised in
Rocq: they are relatively simple, non-technical and
their formalisation would require many times more
preliminary work than seems reasonable.)

As requested further, LLM confessed that it was largely
inspired by studies on complexity of normalisation for
STLC [CITE] [CITE]. And indeed, upon taking a closer
look at the proof of Theorem~\ref{normalisation}, one
cannot help but notice it actually being a primitive
recursive [CITE] algorithm for term normalisation in
disguise. Or something a tad faster:

\begin{definition}[$d$-bounded terms]
    A PTS-$rk$ term is $d$-bounded if it is well-typed
    in some well-formed context and heads of all
    redexes in it have rank equal to or below $d$.
\end{definition}

\begin{theorem}
[Elementary recursive normalisation of bounded
    PTS-$rk$]
\label{bounded-sn}
For each $d$, there is an elementary recursive [CITE]
algorithm for finding a normal form of $d$-bounded
terms.
\end{theorem}
\begin{proof}
Consider the following relation of
``greedy parallel reduction'':
\begin{mathpar}
\inferrule{}{x\Rightarrow x\\ s\Rightarrow s}\and
\inferrule
    {T\Rightarrow T'\\ U\Rightarrow U'}
    {\Pi(x:T).U\Rightarrow\Pi(x:T').U'}\and
\inferrule
    {T\Rightarrow T'\\ t\Rightarrow t'}
    {\lambda(x:T).t\Rightarrow\lambda(x:T').t'}\and
\inferrule
    {t\Rightarrow(\lambda(x:T).t')\\ u\Rightarrow u'}
    {t\,u\Rightarrow t'[u'/x]}\and
\inferrule
    {t\Rightarrow t'\\ t'\ne\lambda\ldots
    \\ u\Rightarrow u'}
    {t\,u\Rightarrow t'\,u'}
\end{mathpar}
It differs from usual parallel reduction since it
greedily applies every ``many-argument'' redex
available in a term. For example,
$(\lambda(x:X).\lambda(y:Y).\lambda(z:Z).t)\,a\,b\,c$
gets reduced to $t'[a/x][b/y][c/z]$ in a single step.
This relation has the following properties:
\begin{itemize}
    \item $(\Rightarrow)\subseteq(\to_\beta^*)$, so
        it is several usual reductions baked into one.
    \item This relation is total and functional.
    \item This relation can be realised with an
        elementary recursive algorithm $\mathcal{A}$:
        traverse the term bottom-up, immediately
        applying every $\beta$-redex which is formed
        when result of a recursive run of algorithm on
        application head is a lambda.
    \item If $t\Rightarrow t'$, then
        $|t'|\le(2+\varepsilon)^{\mathcal{O}(|t|)}$.
    \item If redex heads in $t$ have ranks less than
        \emph{or equal to} $d$ and $t\Rightarrow t'$,
        redex heads in $t'$ have ranks \emph{less than}
        $d$: new redexes appear only after substitution
        of abstractions into application heads, and
        ranks of these abstractions are strictly less
        than the rank of redex heads in which these
        substitutions happened.
\end{itemize}
Because of the last property, it is enough to run
$\mathcal{A}$ for $d$ times in total to normalise a
$d$-bounded term; since input size on $k$-th run is
bounded by $k$-th tower of exponents, the whole
algorithm is elementary recursive.

To make sure that the corresponding function is total,
output $0$ if result of an algorithm contains a
$\beta$-redex.
\end{proof}

Similar result was already proven for STLC [TODO] since
simple types can be classified by ranks as well. And
since rank is easily computable for simple types,
normalisation of arbitrary STLC terms is in
$\mathsf{TOWER}$. A little weaker result holds for our
dependently typed PTS:

\begin{theorem}
    Normalisation of \textbf{closed} well-typed
    PTS-$rk$ terms is in $\mathsf{TOWER}$.
\end{theorem}
\begin{proof}
    By Theorem~\ref{bounded-sn}, it is enough to
    compute some upper bound $B$ on the rank of types
    appearing in abstractions of a term, then run a
    normaliser for $(B+1)$-bounded terms. Consider the
    following elementary recursive functions, written
    in pseudocode and using a suitable encoding of
    terms, tuples and lists:
    \begin{minted}{agda}
maxRankLiteral rkᵣ      = r
maxRankLiteral Π(x:T).U =
    maxRankLiteral T ⊔ maxRankLiteral U
maxRankLiteral x        = 0
maxRankLiteral λ(x:T).t =
    maxRankLiteral T ⊔ maxRankLiteral t
maxRankLiteral (t u)    =
    maxRankLiteral t ⊔ maxRankLiteral u

traversal _ rkᵣ      _ = (r , 0 , 0)
traversal Γ Π(x:T).U d =
    let (vᵀ , rᵀ , aᵀ) = traversal Γ T d
        (vᵁ , rᵁ , aᵁ) = traversal (x := vᵀ, Γ) U d
     in (vᵁ , suc rᵀ ⊔ rᵁ , aᵀ ⊔ aᵁ)
traversal Γ x        d = (d , Γ[x] , 0)
traversal Γ λ(x:T).t d =
    let (vᵀ , rᵀ , aᵀ) = traversal Γ T d
        (vᵗ , rᵗ , aᵗ) = traversal (x := vᵀ, Γ) t d
     in (vᵗ , rᵗ , rᵀ ⊔ aᵀ ⊔ aᵗ)
traversal Γ (t u)    d =
    let (vᵗ , rᵗ , aᵗ) = traversal Γ t d
        (_  , _  , aᵘ) = traversal Γ u d
     in (vᵗ , rᵗ , aᵗ ⊔ aᵘ)

estimateRank t =
    let (vᵗ , rᵗ , aᵗ) =
        traversal [] t (maxRankLiteral t)
     in aᵗ
    \end{minted}
    Function $\mathtt{maxRankLiteral}\;t$ computes
    maximal rank which is contained as a subterm in
    term $t$. The workhorse of this algorithm,
    $\mathtt{traversal}\;\Gamma\;t\;d$, accepts the
    list $\Gamma$ of ranks in context of a term $t$ and
    a ``default rank'' $d$, and outputs the triple
    $(v,r,a)$, where $v$ is value of $t$ as a rank (if
    supplied with enough arguments), $r$ is its rank if
    $t$ is a type (if supplied with enough arguments)
    and $a$ is an upper bound on rank of types
    appearing in abstractions which we are trying to
    compute. The only suspicious case is
    $\mathtt{traversal}\,\Gamma\,x\,d$ where we output
    $d$ (the ``default'' rank) as a value of $x$:
    could we compute value of $x$ instead, if present?
    Regrettably, this would require normalisation of a
    term corresponding to $x$, which is too long;
    instead we notice that $\mathsf{rk}_r$ can be
    reduced only from terms which already contain it as
    a subterm, so value of $x$ as a rank can only be
    $\mathsf{rk}_r$ which is already present somewhere
    in original term. And our result is monotone with
    respect to the value we output here so it is enough
    to provide some upper bound on it; thus it is
    enough to output a maximal rank which occurs in
    original term as a subterm.
\end{proof}

Note that closedness condition is unavoidable: without
it, we would not know ranks of type variables from
context and there would be no known bound to the rank
of redexes.

Another consequence of Theorem~\ref{bounded-sn} is a
natural upper bound for functions \textit{uniformly}
expressible in PTS-$rk$ (this result is also similar to
the one for STLC [CITE]).

\begin{definition}
    A family of functions on natural numbers
    $\mathcal{F}$ is \emph{uniformly} expressible in
    PTS-$rk$ if there is a \emph{natural number type}
    $N:\mathsf{rk}_r$, a \emph{zero} $Z:N$, and a
    \emph{successor} $S:N\to N$ such that for each
    $(f:\mathbb{N}^k\to\mathbb{N})\in\mathcal{F}$ there
    is an \emph{encoding} $F:N\to\ldots\to N\to N$ such
    that
    \[
        \forall (n_1,\ldots,n_k)\in\mathbb{N}^k:
        \quad F\,(S^{n_1}\,Z)\,\ldots\,(S^{n_k}\,Z)
        =_\beta S^{f(n_1,\,\ldots,\,n_k)}\,Z
    \]
\end{definition}

\begin{theorem}\label{uniform-elementary}
    A family of elementary recursive functions
    is not uniformly expressible in PTS-$rk$.
\end{theorem}
\begin{proof}
    By direct diagonalisation argument. Suppose we
    could encode all elementary recursive functions
    with $N:\mathsf{rk}_r$, $Z:N$ and $S:N\to N$.

    Denote by $z$ and $s$ normal forms of $Z$ and $S$,
    respectively. By Theorem~\ref{bounded-sn},
    normalisation of $(r+1)$-bounded terms
    $\mathcal{SN}_{r+1}$ is elementary recursive, so
    the following function is elementary recursive as
    well:
    \[
        h(t) = \begin{cases}
            1,&\mathcal{SN}_{r+1}(t\,(s^{[t]}\,z))=z,\\
            0&\text{otherwise}.
        \end{cases}
    \]
    (Where $[t]\in\mathbb{N}$ is code of a term $t$.)

    By assumption, it is encodable by
    $\mathcal{H}:N\to N$; denote by $H$ its normal
    form. Observe the value of $h(H)$:
    \[
        h(H)=1\Leftrightarrow
        \mathcal{SN}_{r+1}(H\,(s^{[H]}\,z))=z
        \Leftrightarrow
        H\,(s^{[H]}\,z) =_\beta z
        \Leftrightarrow
        h(H)=0
    \]
    First equivalence is by definition of $h$; second
    is by virtue of expression under $\mathcal{SN}$
    being $(r+1)$-bounded; the last is by definition
    of $h$ being encoded by $\mathcal{H}$.
\end{proof}

A logical next question is then: which functions are
computable? Elementary recursive! Though, of course,
not in the same sense as before.

\begin{definition}[Polymorphic numeral family]
    For each fixed $j\ge 0$, there is a type of
    \emph{polymorphic Church numerals} of rank $j+2$:
    \[ \mathcal{N}_j \triangleq \Pi(X:\mathsf{rk}_j).
        (X\to X)\to X\to X;
    \]
    and for each fixed $n$, there is a corresponding
    \emph{Church literal} of type $\mathcal{N}_j$:
    \[
        \overline{n}\triangleq
        \lambda\,X\,f\,x.f^n\,x.
    \]
    ($f^n\,x$ means repeated application of $f$ to $x$:
    $f(\ldots(f\,x)\ldots)$.)
\end{definition}

\begin{theorem}
    [Non-uniform elementary representability]
    For each elementary recursive function
    $f:\mathbb{N}^n\to\mathbb{N}$, there is a
    \emph{cost} $c_f$ such that for each \emph{base
    rank} $j$, there is an \emph{encoding} $F_j:
    \mathcal{N}_{j+c_f}\to\ldots\to\mathcal{N}_{j+c_f}
    \to\mathcal{N}_j$ such that
    \[
        \forall (x_1,\ldots,x_n)\in\mathbb{N}^n:\quad
        F_j\,\overline{x_1}\,\ldots\,\overline{x_n}
        =_\beta\overline{f(x_1,\ldots,x_n)}.
    \]
\end{theorem}
\begin{proof}
    As is well-known, simply typed Church numerals
    admit addition and multiplication [CITE]. It is
    still true for our rank-refined version:
    \[\begin{array}{rcll}
        \mathsf{Add}_j &\triangleq&
        \lambda\,m\,n\,X\,f\,x.
            m\,X\,f\,(n\,X\,f\,x) &:
            \mathcal{N}_j\to\mathcal{N}_j\to
            \mathcal{N}_j \\
        \mathsf{Mul}_j &\triangleq&
        \lambda\,m\,n\,X\,f\,x.
            m\,X\,(n\,X\,f)\,x &:
            \mathcal{N}_j\to\mathcal{N}_j\to
            \mathcal{N}_j
    \end{array}\]
    By substitution-basis theorem [CITE], it is enough
    to additionally provide modulus
    $x\,\mathrm{mod}\,y$ (where $x\,\mathrm{mod}\,0=x$)
    and binary exponentiation $2^x$ since substitution,
    constants and projections are ``free'' in our
    calculus.

    By Theorem~\ref{uniform-elementary}, we cannot hope
    to have both remaining operations represented
    uniformly. Indeed, na\"{i}ve import of traditional
    exponentiation term for Church numerals [CITE]
    raises the rank of an exponent:
    \begin{multline*}
        \mathsf{Exp}_j \triangleq
        \lambda\,m\,n\,(X:\mathsf{rk}_j)\,f\,x.\,
        n\,(X\to X)\,(\lambda (h:X\to X)(y:X).
            m\,X\,h\,y
        )\,f\,x
        \\:\mathcal{N}_j\to\mathcal{N}_{j+1}\to
        \mathcal{N}_j
    \end{multline*}
    And that is exactly the reason we allowed input and
    output ranks of Church numerals to differ.

    There still remains a problem, however: by allowing
    ranks to differ, we lost uniformity of
    substitution. In particular, consider an elementary
    function $x\mapsto x + 2^x$. To represent such
    computation using combinators we defined, $x$ has
    to have both type $\mathcal{N}_j$ (as left addend)
    and $\mathcal{N}_{j+1}$ (as exponentiation power).
    Instead of introducing subtyping in the system, we
    make use of the following family of
    \emph{numeral-preserving lowerings}:
    \[\delta_j\triangleq\lambda\,n\,X\,f\,x.\,
        n\,(X\to X)\,(\lambda\,h\,y.f\,(h\,y))\,
        (\lambda\,y.y)\,x
        :\mathcal{N}_{j+1}\to\mathcal{N}_j
    \]
    As the name implies, the goal of such combinators
    is to lower the rank of a numeral, but preserve its
    ``value''. It is achieved in the following way:
    given an $f$ and $x$ for lower-ranked type, an
    original combinator for $n$ is used to build an
    $n$-fold composition of $f$ with itself which is
    then applied to $x$.

    Thus we can lower ranks arbitrarily: $\delta_{b,a}
    \triangleq\delta_a\circ\ldots\circ\delta_{b-1}:
    \mathcal{N}_b\to\mathcal{N}_a$. This restores
    substitution and allows gluing representation of
    elementary recursive functions together; cost $c_f$
    of a function can be computed from any expression
    that represents it --- addition and multiplication
    cost $0$, exponentiation costs $1$, operands are
    merged by taking the maximum of their costs. (Of
    course, cost is not unique and depends on the way
    function is represented with an expression.)

    All that is left, then, is to represent modulus.
    General idea is as follows: first, suppose that
    we are dealing with a simpler problem of taking
    modulus by a constant modulo, e.g.
    $\mathrm{mod}\,5$. Notice that there are only $5$
    possible answers: $0$,$1$,$2$,$3$,$4$. As such,
    they can be represented by the type which contains
    exactly $5$ elements; Church encoding [CITE] yields
    a type of \emph{selectors}:
    \[V_5 \triangleq \Pi(X:\mathsf{rk}_j).\,
        X\to X\to X\to X\to X\to X
        :\mathsf{rk}_{j+1}
    \]
    Its $5$ different canonical values simply pick one
    of $5$ inputs, to each their own.

    There is an operation of \emph{cyclic shift} on
    selectors:
    \[c_5\triangleq \lambda\,s\,X\,x_1\,x_2\,x_3\,x_4
    \,x_5.\,s\,X\,x_2\,x_3\,x_4\,x_5\,x_1
    :V_5\to V_5\]
    Iterating it $n$ times on some pinned ``zeroth''
    selector, we actually get the
    $(n\,\mathrm{mod}\,5)$-th selector. All that is
    left is to read the numeral back:
    \[\rho\triangleq \lambda\,s\,X\,f\,x.\,s\,X\,
    (f^0\,x)\,(f^1\,x)\,(f^2\,x)\,(f^3\,x)\,(f^4\,x):
        V_5\to\mathcal{N}_j\]
    Now, in order to model $m\,\mathrm{mod}\,n$ in this
    vein, we need to \emph{generate} $V_n$ dynamically.
    Lucky for us, PTS-$rk$ is dependently typed so
    there is hope.

    And indeed, the following construction was produced
    in discussion with GPT-6 Astra.

    We are going to produce a combinator $\mathsf{Mod}:
    \mathcal{N}_{j+1}\to\mathcal{N}_{j+4}\to
    \mathcal{N}_j$. Fix parameter $X:\mathsf{rk}_j$ and
    consider the following signature for abstract
    $S:\mathsf{rk}_{j+1}$:
    \[\begin{array}{ll}
        K:X\to S, &z:S,\\
        L:(X\to S)\to X\to S, &c:S\to S,\\
        \rho:S\to(X\to X)\to X\to X,
        &r:(X\to X)\to(X\to X).
    \end{array}\]
    Intention is that $S$ is simply-typed $V_q$ for
    $X$, $z$ is its ``zero'', $K$ is a ``zero'' for
    $V_{q+1}$, $c$ is a cyclic shift for $V_q$, $L$ is
    a cyclic shift for $V_{q+1}$, $\rho$ converts a
    selector into a numeral, and $r$ is the result of
    modulus computed using structure from the previous
    step (this is similar to the trick used to define a
    predecessor function).

    A callback accepting such structure and returning
    $X$ has the following type:
    \begin{align*}
        H_X \triangleq\;&\Pi(S:\mathsf{rk}_{j+1}).\;
        (X\to S)\to S\to((X\to S)\to X\to S)\to(S\to S)
        \to\\&(S\to (X\to X)\to X\to X)\to
        ((X\to X)\to X\to X)\to X
        \quad :\mathsf{rk}_{j+3}
    \end{align*}
    A type of ``records'' is thus
    $P_X\triangleq H_X\to X:\mathsf{rk}_{j+4}$, with a
    constructor
    \[
        \langle S,K,z,L,c,\rho,r\rangle\triangleq
        \lambda h.\,h\,S\,K\,z\,L\,c\,\rho\,r.
    \]
    As an initial $P_X$, we build a structure for
    $V_1$ and an $r$ which behaves like
    $m\,\mathrm{mod}\,0=m$:
    \begin{align*}
        S_0&\triangleq X\to X,\\
        K_0&\triangleq\lambda(a:X)(b:X).a,&
        z_0&\triangleq\lambda(a:X).a,\\
        L_0&\triangleq\lambda(t:X\to S_0)(a:X)(b:X).
            t\,b\,a,&
        c_0&\triangleq\lambda(s:S_0).s,\\
        \rho_0&\triangleq
            \lambda(s:S_0)(g:X\to X)(a:X).s\,a,&
        r_0&\triangleq\delta_j\,m\,X.
    \end{align*}
    Pack them into $p_0\triangleq\langle\ldots\rangle$.

    Then, given a structure for $q$, a structure for
    $q+1$ is:
    \begin{align*}
        S^+&\triangleq X\to S,\\
        K^+&\triangleq\lambda(a:X)(b:X).Ka,&
        z^+&\triangleq K,\\
        L^+&\triangleq\lambda(t:X\to S^+)(a:X)(b:X).
            L\,(t\,b)\,a,&
        c^+&\triangleq L,\\
        \rho^+&\triangleq\lambda(t:S^+)(g:X\to X)(a:X).
            \rho\,(ta)\,g\,(ga),&
        r^+&:=\rho\,(m\,S\,c\,z).
    \end{align*}
    A transition term can be written like this:
    \[
        T\triangleq\lambda(p:P_X)(h:H_X).
            p\,(\lambda\,S\,K\,z\,L\,c\,\rho\,r.\,
            h\,S^+\,K^+\,z^+\,L^+\,c^+\,\rho^+\,r^+)
        :P_X\to P_X
    \]
    Therefore, the following term computes modulus with
    expected behaviour:
    \begin{multline*}
        \mathsf{Mod}_j\triangleq\lambda
        (m:\mathcal{N}_{j+1})(n:\mathcal{N}_{j+4})
        (X:\mathsf{rk}_j)(f:X\to X)(x:X).\\
            n\,P_X\,T\,p_0\,
            (\lambda\,\_\,\_\,\_\,\_\,\_\,\_\,r.
                r\,X\,f\,x).
    \end{multline*}
\end{proof}

Two questions remain open: a) whether there are
non-elementary functions which are non-uniformly
computable in PTS-$rk$ and b) whether there are any
non-polynomial functions which are uniformly
computable in PTS-$rk$.

\section{Set-theoretical model of PTS-$rk$}
\label{rk-model}

Now, studying computational properties of a calculus is
great, but a logic is only useful when it can tell
right from wrong. In this section, we are going to
provide a nontrivial set-theoretical semantics for
PTS-$rk$ with respect to which it is sound, and where
$\llbracket\Pi(X:\mathsf{rk}_j).X\rrbracket$ is going
to be an empty set, as expected.

We are going to need an unusual assumption however:

\begin{theorem}
    PTS-$rk$ has a sound model in ZF+GCH.
\end{theorem}
\begin{proof}
    Produced in discussion with GPT 6 Astra.

    Take distinct codes \(c_\kappa\) for every cardinal
    below \(\aleph_\omega\) and distinct additional
    codes \(s_j\) for \(j\in\mathbb N\). Define
    \[
        U_i=\{c_\kappa:\kappa\le\aleph_i\}\cup
            \{s_j:j\in\mathbb N\},
        \qquad
        \operatorname{El}(c_\kappa)=\kappa,\quad
        \operatorname{El}(s_j)=U_j.
    \]
    That is,
    \begin{align*}
        U_0 &= \{c_0,c_1,\ldots,c_\omega=c_{\aleph_0},
                s_0,s_1,\ldots\},\\
        U_1 &= \{c_0,c_1,\ldots,c_{\aleph_0},
                c_{\aleph_1},s_0,s_1,\ldots\},\\
        \ldots\\
        U_k &= \{
            c_0,c_1,\ldots,c_{\aleph_0},
            c_{\aleph_1},\ldots,c_{\aleph_k},
            s_0,s_1,\ldots\}\\
        \ldots
    \end{align*}
    (Note that $x\in U_i\Rightarrow|\mathrm{El}(x)|\le
    \aleph_i$ since all $U_i$ are countable.)

    Interpretation then goes as follows:
    \begin{itemize}
        \item Interpret \(\mathsf{rk}_i\) by \(s_i\).
            The sort axiom holds since
            \(s_i\in U_0=\operatorname{El}(s_0)\).
        \item If \(a\in U_m\) and \(b(x)\in U_n\) for
            every \(x\in\operatorname{El}(a)\), the
            ordinary dependent function set has
            cardinality at most
            \[
                \left|\prod_{x\in\operatorname{El}(a)}
                \operatorname{El}(b(x))\right|
                \le \aleph_n^{\aleph_m}
                \le \aleph_{\max(m+1,n)}.
            \]
            (The last inequality follows from GCH since
            alephs have finite indices.)

            Thus, for all such $a,b$ we can choose a
            cardinal $\kappa_{\Pi(x:a).b(x)}$ (whose
            code lies in $U_{\max(m+1,n)}$) along with
            a bijection
            \[
                f:\left(\prod_{
                    x\in\operatorname{El}(a)}
                \operatorname{El}(b(x))
                \right)\simeq\kappa_{\Pi(x:a).b(x)}
            \]
            The eligible pairs $(\kappa, f)$ form a
            set, so ordinary choice (which holds in
            ZF+GCH [CITE]) suffices. Make each choice
            once for each $(a,b)$, independently of a
            typing derivation or rank bound.
        \item Abstraction and application are modeled
            by $f$ and $f^{-1}$ respectively, as
            defined above.
        \item Contexts are interpreted by valuations.
    \end{itemize}
    It is easy to check that typing rules are respected
    by such interpretation, making it a sound model.
\end{proof}

A few corollaries:

\begin{itemize}
    \item Indeed, the product interpreting
        \(\Pi Z:\mathsf{rk}_0.Z\) ranges over \(U_0\),
        which contains \(c_0\), whose decoding is
        $\varnothing$. That product is therefore empty.
    \item A system with rank cumulativity
        $t:\mathrm{rk}_j\Rightarrow
        t:\mathrm{rk}_{j+1}$ is sound with respect to
        this model; normalisation of bounded terms
        remains elementary as in
        Theorem~\ref{bounded-sn}.
\end{itemize}

Note that usage of GCH is unavoidable in this model to
validate the rules. We could try building a model in
ZFC using $\beth$-s instead of $\aleph$-s, but that
would break an assumption that $U_i$ are countable,
which would falsify axioms
$\mathsf{rk}_n:\mathsf{rk}_0$. That is not a huge
problem however since ZF+GCH is interpretable in ZFC
[CITE] (in contrast to ZFC + inaccessible cardinals
which is not [CITE]).

\section{Proof-theoretic strength of PTS-$rk$}

To demonstrate ability of PTS-$rk$ to embed nontrivial
logical theories, we construct a fixed interpretation
of Tarski--Mostowski--Robinson arithmetic $\R$ in
$\PTS$. A unary Church domain and certified finite
observations first interpret its relational form $\Rt$.
Definable graphs then yield $\R$; a direct
interpretation with arithmetic operators is also given.
All constructions use only dependent function types and
$\beta$-conversion.

\subsection
{Classical arithmetic and the target calculus}
\label{sec:logic}

We work in classical first-order logic with equality.
Let $\Rt$ be the relational presentation of
Tarski--Mostowski--Robinson arithmetic in the language
$\{0,S,=,\le,A,M\}$, with ternary predicates $A,M$ and
numerals $\bar n=S^n0$. Its axioms are the universal
closures of the following schemes, indexed by external
$m,n\in\mathbb N$:
\begin{align}
 A(\bar m,\bar n,x)&\leftrightarrow \overline{m+n}=x,
    \tag{$\widetilde R1$}\label{rtilde:add}\\
 M(\bar m,\bar n,x)&\leftrightarrow \overline{mn}=x,
    \tag{$\widetilde R2$}\label{rtilde:mul}\\
 \bar m&\ne\bar n\quad(m\ne n),
    \tag{$\widetilde R3$}\label{rtilde:neq}\\
 x\le\bar n&\leftrightarrow\bigvee_{i=0}^{n}x=\bar i,
    \tag{$\widetilde R4$}\label{rtilde:bound}\\
 x\le\bar n&\lor\bar n\le x.
    \tag{$\widetilde R5$}\label{rtilde:compare}
\end{align}
Here $m+n$ and $mn$ are calculated externally: the
language has no binary arithmetic function symbols.
Section~\ref{sec:nonrelational} relates this theory
to the original, functional arithmetic
$\R$~\cite{tarski1953undecidable}.

Put
\[
 \bot:=((Z:\rk0)\to Z):\rk1,\qquad
 \top:=\rk0\to\rk0:\rk1,\qquad
 \mathsf{id}:=\lambda(Z:\rk0).Z:\top.
\]
Use the fixed-answer proof types
\begin{gather*}
 \nc P:=P\to\bot,\qquad
 P\ac Q:=(P\to Q\to\bot)\to\bot,\\
 \bigvee_c(P_0,\ldots,P_k):=
 (P_0\to\bot)\to\cdots\to(P_k\to\bot)\to\bot,\\
 P\iffc Q:=(P\to Q)\ac(Q\to P).
\end{gather*}
A type $P$ is \emph{stable} when $\nc\nc P\to P$ is
inhabited. Negative types, dependent functions with
stable codomain, and these encoded conjunctions and
disjunctions are stable. They support the classical
logical rules on stable proof types; case elimination
is into stable targets. In particular, $P\oc\nc P$ has
the proof
$\lambda(k:P\to\bot)(l:(P\to\bot)\to\bot).\,l\,k$.
We never eliminate $\bot$ into an arbitrary positive
type.

For a carrier $V$, a domain predicate $\Delta$, and
fixed stable atomic interpretations, use implication as
an arrow and translate quantifiers by
\begin{align*}
 (\forall x\,\varphi)^*&:=(x:V)\to\Delta(x)\to\varphi^*(x),\\
 (\exists x\,\varphi)^*&:=
 ((x:V)\to\Delta(x)\to\varphi^*(x)\to\bot)\to\bot.
\end{align*}
All translated formulas are stable. Once domain closure
and atomic congruence are proved, the logical rules and
equality substitution extend structurally to formulas
on admitted arguments. All interpretations below use
this one logical encoding. For the relational
interpretation, Section~\ref{sec:relational}
establishes completeness of the ground proof/refutation
construction and failure of faithfulness on arbitrary
first-order sentences.

For exact-rank motives, let $J_1=\top$,
$J_{k+1}=J_k\to J_k$, with $j_1=\mathsf{id}$ and
$j_{k+1}=\lambda(x:J_k).x$. If $P:\rk a$, $a\le k$, and
$k\ge2$, then
\[
 L_k(P):=J_{k-1}\to P:\rk k.
\]
Abstraction and evaluation at $j_{k-1}$ give maps
between $P$ and $L_k(P)$.

\subsection{Unary domain and equality}\label{sec:unary}
Use the Church numeral carrier
\[
 D:=\bigl((X:\rk1)\to X\to(X\to X)\to X\bigr):\rk3,
 \qquad
 \z:=\lambda X\,z\,s.z,\qquad
 \suc d:=\lambda X\,z\,s.s(d\,X\,z\,s).
\]
Write $\nu_n=\suc^n\z$. Define the two auxiliary,
same-type Church operations
\begin{align*}
 \Add\,d\,e&:=\lambda X\,z\,s.\ d\,X\,(e\,X\,z\,s)\,s,
    \\
 \Mul\,d\,e&:=\lambda X\,z\,s.\ d\,X\,z\,(\lambda y:X.
    e\,X\,y\,s).
\end{align*}
Both have type $D\to D\to D$. On canonical numerals
they compute addition and multiplication by beta
reduction. No eta equation on arbitrary $d:D$ is used.

The domain predicate is the following
constructor-closure certificate:
\[
 \Good(d):=\bigl((Q:D\to\rk{10})\to Q(\z)
 \to((u:D)\to Q(u)\to Q(\suc u))\to Q(d)\bigr):\rk{12}.
\]
Zero and successor preserve it by abstraction. For
example,
\[
 \mathsf{good}_0:=\lambda Q\,q_0\,q_s.\,q_0,\qquad
 \mathsf{good}_s\,d\,g:=\lambda Q\,q_0\,q_s.\,
 q_s\,d\,(g\,Q\,q_0\,q_s).
\]
Thus the interpreted domain is nonempty and closed
under the source functions. The auxiliary $\Add,\Mul$
occur only inside predicates, so no closure of $\Good$
under them is required.

Define equality by negative tests:
\[
 \boxed{E(d,e):=\bigl((P:D\to\rk1)\to
       ((\nc P(d))\iffc(\nc P(e)))\bigr):\rk5.}
\]
Its reflexivity, symmetry, and transitivity follow
pointwise from the corresponding rules for logical
equivalence. Every $F:D\to D$ respects $E$: to obtain
$E(Fd,Fe)$ from $h:E(d,e)$, specialize $h$ to $u\mapsto
P(Fu)$. This applies to successor and to each argument
of the auxiliary arithmetic operations. Equality
substitution for $E$ itself follows from symmetry and
transitivity. Define the arithmetic atoms by
\[
 A^*(d,e,f):=E(\Add\,d\,e,f),\qquad
 M^*(d,e,f):=E(\Mul\,d\,e,f).
\]
They respect $E$ in all arguments. Their ground-input
defining schemes are immediate from the beta equations,
including an arbitrary admitted third argument.

\subsection{Order from successor-closed unary tests}
\label{sec:order}
A certified unary algebra is the fixed telescope
\[
 \begin{aligned}
 \mathcal U={}&(X:\rk1)\,(z:X)\,(s:X\to X)\,
     (\delta:X\to\rk4)\\
             &(d_0:\delta z)\,
      (d_s:(u:X)\to\delta u\to\delta(su)).
 \end{aligned}
\]
Writing $\mathcal U\to B$ abbreviates the successive
dependent arrows for these binders; it does not
introduce a record type. Put $\ev_{\mathcal U}(d)=d\,X
\,z\,s$ and $G(d):=(\mathcal U\to\delta(\ev_{\mathcal
U}(d))):\rk6$. For $P:X\to\rk1$, write
\[
 \Up_{\mathcal U}(P):=
 (u:X)\to\delta u\to\nc P(u)\to\nc P(su).
\]
Define
\[
 \boxed{\begin{aligned}
 O(d,e):=\bigl(\mathcal U\to(P:X\to\rk1)\to{}&
     \Up_{\mathcal U}(P)\to{}\\
 &\nc P(\ev_{\mathcal U}(d))\to\nc P(\ev_{\mathcal U}
     (e))\bigr):\rk6.
 \end{aligned}}
\]
This needs only $d_0,d_s$, and an upward-closure
certificate. Preorder laws and equality compatibility
are \emph{derived}, not included as certificates.

\begin{lemma}\label{basic}
    $G(\z)$ and $G(d)\to G(\suc d)$ are inhabited, as
    is $\Good(d)\to G(d)$. The relation $O$ is
    reflexive and transitive on all raw $D$-objects.
    Moreover, $E(d,e)\to O(d,e)$ and $G(d)\to O(d,\suc
    d)$ are inhabited. Both arguments of $O$ respect
    $E$.
\end{lemma}
\begin{proof}
    The two $G$-closure terms use $d_0,d_s$ and the
    beta fold equations. To obtain $G(d)$ from
    $\Good(d)$, instantiate its predicate at
    $u\mapsto L_{10}(G(u))$ and evaluate the result at
    $j_9$.

    Reflexivity and transitivity of $O$ are identity
    and composition on each test. Given $E(d,e)$,
    instantiate its predicate with $u\mapsto P(
    \ev_{\mathcal U}(u))$ and take the forward
    implication. The step law is the upward-closure
    certificate applied to $\ev_{\mathcal U}(d)$ and
    the supplied $G$-evidence. Finally use
    $E\Rightarrow O$, symmetry of $E$, and transitivity
    of $O$ for compatibility.
\end{proof}
Consequently, for each external $i\le j$, there is a
closed proof of $O(\nu_i,\nu_j)$, built by a finite
succession of steps.

\subsection{Finite separation}\label{sec:separation}
\begin{lemma}\label{sep}
    For every external $n$, there is a closed term
    $\mathsf{no\_back}_n:O(\nu_{n+1},\nu_n)\to\bot$.
    Consequently distinct standard numerals are
    internally $E$-unequal, and all reversed ground
    $O$-comparisons have internal refutations.
\end{lemma}
\begin{proof}
    Put $c=n+1$, $q=c+1$, and use the rank-one
    selectors
    \[
        F_q:=\underbrace{\rk1\to\cdots\to\rk1}_{
            q\ \mathrm{arguments}}
        \to\rk1,\qquad
        e_i:=\lambda Z_0\ldots Z_{q-1}.Z_i.
    \]
    Choose $X=F_q$, $z=e_0$, and the unary table
    \[
        s(u):=\lambda Z_0\ldots Z_{q-1}.\
        u\,Z_{\min(0+1,c)}\cdots Z_{\min(c+1,c)}.
    \]
    Thus $s(e_i)=_\beta e_{\min(i+1,c)}$. Let
    \[
        \delta(u):=\bigl((Q:F_q\to\rk3)\to
        Q(e_0)\to\cdots\to Q(e_c)\to Q(u)\bigr):\rk4.
    \]
    Selector projection gives $d_0$. To construct
    $d_s$, open the output $\delta(su)$ at $Q$ and its
    branch proofs, and specialize the input certificate
    to $v\mapsto Q(sv)$. Its branches are the supplied
    proofs for the appropriate successor selectors.

    Use the unary test
    \[
        P(u):=u\,\underbrace{\top\cdots\top}_{
            c\ \mathrm{entries}}\,\bot.
    \]
    To prove $\Up(P)$, eliminate $\delta(u)$ into the
    rank-three motive $v\mapsto\nc P(v)\to\nc P(sv)$.
    At $e_c$ both tests are $\bot$ and the implication
    is identity. At $e_i$ with $i<c$, a premise
    $\top\to\bot$ gives $\bot$ by application to
    $\mathsf{id}:\top$; abstraction gives the negative
    conclusion.

    Specializing $h:O(\nu_c,\nu_n)$ to these data
    yields $(\bot\to\bot)\to(\top\to\bot)$. Apply it to
    $\mathsf{id}_{\bot}$ and $\mathsf{id}_{\top}$.

    If $m>n$, compose $O(\nu_{n+1},\nu_m)$ with a
    purported $O(\nu_m,\nu_n)$ and use
    $\mathsf{no\_back}_n$. A purported $E(\nu_m,\nu_n)$
    supplies this last order proof by
    Lemma~\ref{basic}; symmetry handles $m<n$.
\end{proof}
The selectors are certified only through $\delta$; no
parametricity claim about all raw inhabitants of $F_q$
is needed. These are finite separating observations,
not finite models of the entire source theory.

\subsection{Bounded cases}\label{sec:cases}
Set
\[
 \Case_n(d):=\bigvee_c\bigl(E(d,\nu_0),\ldots,
 E(d,\nu_n),O(\nu_{n+1},d)\bigr)
       :\rk8.
\]
\begin{lemma}\label{case}
    $(d:D)\to\Good(d)\to\Case_n(d)$ is inhabited for
    every external $n$.
\end{lemma}
\begin{proof}
    Use the rank-ten positive invariant
    \[
        \Inv_n(d):=(Z:\rk8)\to(G(d)\to\Case_n(d)\to Z)
        \to Z.
    \]
    At zero, supply $G(\z)$ and the $E(\z,\nu_0)$
    branch. At successor, open the target continuation
    and eliminate the child invariant at that same $Z$,
    obtaining $G(d)$ and $\Case_n(d)$. Obtain
    $G(\suc d)$ independently.
    To construct the new negative case proof, all
    branch continuations return $\bot$. An $i<n$ branch
    becomes $i+1$ by successor congruence. The $i=n$
    branch gives $E(\suc d,\nu_{n+1})$, hence
    $O(\nu_{n+1},\suc d)$. The large branch is composed
    with $O(d,\suc d)$. This proves successor closure.

    Instantiate $\Good(d)$ with $\Inv_n$. Project its
    case component by choosing $Z=\Case_n(d)$ and
    continuation $\lambda g\,c.c$. No negative
    disjunction is eliminated into positive
    $G$-evidence.
\end{proof}

\subsection{The relational interpretation}
\label{sec:relational}

\begin{theorem}\label{main}
    The carrier $D$, admitted domain $\Good$, closed
    operations $\z,\suc$, and relations $E,O,A^*,M^*$
    interpret the classical relational theory $\Rt$ in
    $\PTS$. Every translated axiom has an effectively
    constructible closed proof term. The interpretation
    is closed under classical first-order derivations:
    every $\Rt$-derivation of a sentence effectively
    yields a closed $\PTS$ proof term of its
    translation. True ground arithmetic atoms have
    proofs, and false ground atoms have internal
    refutations. The interpretation is not, however,
    faithful on arbitrary first-order sentences.
\end{theorem}
\begin{proof}
    Use Section~\ref{sec:logic}'s translation with
    $V=D$ and $\Delta=\Good$. Domain closure, atomic
    stability, and congruence have been established;
    hence source deductions translate to PTS proofs.

    For the first two arithmetic schemes,
    $\Add\,\nu_m\,\nu_n=_\beta\nu_{m+n}$ and
    $\Mul\,\nu_m\,\nu_n=_\beta\nu_{mn}$, so both sides
    of each translated biconditional are
    beta-convertible. Numeral inequality is
    Lemma~\ref{sep}.

    For bounded order, apply Lemma~\ref{case}. Under
    $O(d,\nu_n)$ its large branch contradicts
    $\mathsf{no\_back}_n$ by transitivity. Conversely,
    an $E(d,\nu_i)$ branch with $i\le n$ gives
    $O(d,\nu_n)$ using $E\Rightarrow O$ and the ground
    comparison. For comparability the finite branches
    give $O(d,\nu_n)$ and the large branch gives
    $O(\nu_n,d)$.

    The ground proof/refutation claims follow from the
    arithmetic beta equations, reflexivity, ground
    order, and Lemma~\ref{sep}.
\end{proof}
\begin{remark}
    The fixed ranks are $D:3$, $E,A^*,M^*:5$, $G,O:6$,
    $\delta:4$, $\Case:8$, $\Inv:10$, and $\Good:12$.
    The proof uses only $d_0,d_s$ and $\Up(P)$ as
    observation certificates.
\end{remark}

\paragraph{Completeness on ground arithmetic.}
Every source axiom instance has a proof, and every
$\Rt$-derivation translates using the same fixed
definitions. For ground arithmetic atoms, the
construction provides proofs of all true instances and
internal refutations of all false instances. This
ground completeness is not reflection of arbitrary
source derivability.

\begin{remark}[Failure of faithfulness]
\label{rem:relational-nonfaithful}
    The interpretation of $\Rt$ is not faithful: the
    translation of $\rho:=\forall x\,(x\le x)$ is
    provable in $\PTS$, although $\Rt\nvdash\rho$.

    \emph{An explicit proof of order reflexivity.}
    Its translation is
    $\rho^*=(d:D)\to\Good(d)\to O(d,d)$. Expanding
    every binder in the observation telescope gives the
    following closed term:
    \[
        \begin{aligned}
            \mathsf{orderRefl}:=
            {}&\lambda(d:D)(g:\Good(d))\,\\
            &(X:\rk1)(z:X)(s:X\to X)
            (\delta:X\to\rk4)\,\\
            &(d_0:\delta z)(d_s:(u:X)\to\delta u
            \to\delta(su))\,\\
            &(P:X\to\rk1)\,\\
            &(p_{\uparrow}:(u:X)\to\delta u\to
                (P(u)\to\bot)\to(P(su)\to\bot))\,\\
            &(h:P(d\,X\,z\,s)\to\bot).\;h.
        \end{aligned}
    \]
    Thus
    \begin{equation}
        \mathsf{orderRefl}:(d:D)\to\Good(d)\to O(d,d).
        \label{eq:order-reflexivity}
    \end{equation}
    Indeed, after the binders are opened, the remaining
    goal is $P(d\,X\,z\,s)\to\bot$, exactly the type of
    $h$. The term simply returns $h$; it uses neither
    the domain proof $g$ nor the three observation
    certificates $d_0,d_s,p_{\uparrow}$. Removing the
    binder $g$ also gives $(d:D)\to O(d,d)$, so the
    proof works on every raw presentation.

    \emph{A relational countermodel.}
    Take $N^+=\mathbb N\sqcup\{\omega\}$, where
    $\omega$ is one new element. Interpret $0$ and
    successor as usual on $\mathbb N$, put
    $S\omega=\omega$, and use actual equality. For
    $x,y\in\mathbb N$, let $A(x,y,z)$ hold exactly when
    $z=x+y$, and let $M(x,y,z)$ hold exactly when
    $z=xy$. When at least one of $x,y$ is $\omega$, let
    each predicate hold exactly when $z=0$. Define
    \[
        x\le^+ y\quad\Longleftrightarrow\quad
        \bigl(x,y\in\mathbb N\ \text{and}
        \ x\le_{\mathbb N}y\bigr)
        \ \text{or}\ \bigl(x\in\mathbb N\ \text{and}
        \ y=\omega\bigr).
    \]
    The ground graph schemes hold for every third
    argument, and distinct numerals are unequal. For
    every standard $n$, the downset of $n$ is exactly
    $\{0,\ldots,n\}$, proving both directions of
    \eqref{rtilde:bound}. Every standard number is
    comparable with $n$, and $n\le^+\omega$, proving
    \eqref{rtilde:compare}. Hence this is a model of
    $\Rt$.

    However, $\omega\not\le^+\omega$, so $\rho$ is
    false. By soundness of classical first-order logic,
    \[
        \Rt\nvdash\rho,
        \qquad\text{whereas}\qquad \PTS\vdash
        \mathsf{orderRefl}:\rho^*.
    \]
    Thus the relational interpretation preserves
    derivability but does not reflect it. The
    countermodel is a model of the source theory, not a
    set-theoretic model of $\PTS$. There is no conflict
    with ground completeness: $\rho$ is universally
    quantified, and its failure occurs at a non-numeral
    element.
\end{remark}

\subsection{Relation to the original presentation of
\texorpdfstring{$\R$}{R}}\label{sec:nonrelational}
The original Tarski--Mostowski--Robinson arithmetic
$\R$~\cite{tarski1953undecidable} uses $\{0,S,+,\cdot,\le\}$ and the
following schemes, again with universal closure:
\begin{align}
 \bar m+\bar n&=\overline{m+n}, &
 \bar m\cdot\bar n&=\overline{mn},
     \tag{R1--R2}\label{r:arithmetic}\\
 \bar m&\ne\bar n\quad(m\ne n), &&
     \tag{R3}\label{r:neq}\\
 x\le\bar n&\longrightarrow\bigvee_{i=0}^{n}x=\bar i, &&
     \tag{R4}\label{r:bound}\\
 x\le\bar n&\lor\bar n\le x. &&
     \tag{R5}\label{r:compare}
\end{align}
Requiring $A,M$ in $\Rt$ to be total, single-valued graphs makes them
interchangeable with the function symbols $+,\cdot$. This functional
strengthening is definitionally equivalent to $\R$. Even without that
strengthening, the two theories have the same interpretability degree.

\begin{proposition}\label{prop:graphs}
The classical theories $\Rt$ and $\R$ are mutually interpretable.
Consequently $\PTS$ interprets $\R$ by definable arithmetic graphs.
\end{proposition}
\begin{proof}
In $\R$, interpret $A(x,y,z)$ by $x+y=z$ and $M(x,y,z)$ by $xy=z$.
The ground arithmetic equations give the two graph schemes.
For $i<j$, an assumption $\bar j\le\bar i$ contradicts
\eqref{r:bound} and numeral separation; \eqref{r:compare} therefore yields
$\bar i\le\bar j$. The case $i=j$ follows from \eqref{r:compare} directly.
Thus $\R$ proves all true ground comparisons, and equality substitution
supplies the reverse implication in \eqref{rtilde:bound}.
The remaining relational schemes are immediate.

For the other direction define, in $\Rt$,
\begin{align*}
 U_A(x,y,z)&:=A(x,y,z)\land
                \forall w\bigl(A(x,y,w)\to w=z\bigr),\\
 F_A(x,y,z)&:=U_A(x,y,z)\lor
       \bigl(z=0\land\neg\exists v\,U_A(x,y,v)\bigr).
\end{align*}
For fixed $x,y$, either a unique $A$-output exists or it does not. In the
first case $F_A$ selects that output; in the second it selects zero.
Hence $F_A$ is provably total and single-valued. At standard inputs,
\eqref{rtilde:add} forces its value to be $\overline{m+n}$.
Define $F_M$ identically with $M$ in place of $A$; its standard value is
$\overline{mn}$. These graphs interpret $+,\cdot$; keep $0,S,=,\le$
unchanged. All five schemes of $\R$ follow.

Composing this graph interpretation with Theorem~\ref{main} gives the
claimed interpretation in $\PTS$. The composition uses the encoded
classical quantifiers to express values and uniqueness; it does not
assert $\Good$-closure of the raw Church operators $\Add,\Mul$.
\end{proof}
The conclusion is mutual interpretability, not bi-interpretability.
In particular, the totalized graphs need not coincide with the original
relations $A,M$ on arbitrary arguments.

\section{Conclusion}
\label{conclusion}

\printbibliography
\end{document}
